\documentclass[../{{cookiecutter.project_slug}}.tex]{subfiles}

\begin{document}

\part{Guidelines}
\chapter{Introduction}
\label{chap:Introduction}

    It is probably a  good idea to split your thesis up  into several files, so
    you can  work on  them individually.  This  is an example  of such  a file,
    which ends up  in the main file  but can be compiled by  itself without all
    the other chapters.  This will save you tons of time.

\chapter{Language}

    \section{I or We?}

    Use the ``we'' form, even when  you're writing alone.  Imagine it to be
    the reader and  you, who are going through the  text together.  It gets
    you closer to your audience.

    \section{Quotes and Emphasis}
    Note that double "quotes" on  your keyboard are not typeset correctly. 
    English  quotes should  look like  ``this'', while  German quotes  are 
    typeset  like \glqq this\grqq.   Note the  difference.  Some  editors 
    insert the correct quotes automatically for you.                       

    Use  emphasis  sparingly.  It  clutters  the  text.  In  general,  use 
    \emph{italics}, which is the one standing out the least.  In captions, 
    it is sometimes useful to use  \textbf{bold} for the signal words left 
    and right, top and bottom: There, you want them to pop out.

    \section{Line Noise and Colloquial Expressions}

    Do  not use  fill words  which do  not carry  meaning.  Try  to be  as 
    specific as  possible.  This does  not mean  as short as  possible, it 
    means that you should have something to say when you write something.
    If you don't know what you want to say, think again, \emph{then} write.

    Don't use colloquial expressions.  In particular, don't use any of the 
    short forms  \emph{don't, aren't,  isn't, you're, we're,  \ldots}.  Do 
    not\footnote{this is much better!} use  ellipsis (\ldots) in the text, 
    as it looks like you're trailing off.                                  

    \chapter{Comments on Writing in LaTeX}

    \section{\LaTeX{}'s Wide Margins}

    Many students  complain  about  large  \LaTeX{}  standard  margins. 
    Typographers have a  rule, though: You shouldn't have  more than about 
    70 characters  per line (when  using single spacing).   Otherwise your 
    eyes have  trouble jumping to  the beginning  of the next  line.  This 
    limits  your column  width severely,  unless you  use a  big font.   A 
    thin  column width  results in  large  margins.  This  is (aside  from 
    portability) the reason  why books are smaller than A4  paper, and why 
    newspapers have multiple thin columns.                                 

    If  you want  to toy  with margins  nevertheless, consider  this: When 
    you  open  your  double-sided  printed thesis,  the  white  region  in 
    the  left/right  and  center  region should  have  equal  widths.   It 
    follows  that outside  margins  are larger  than  the margins  inside. 
    Consequentially, on  a left page,  left margins are larger  than right 
    (i.e.\ center) margins, and on a  right page, right margins are larger 
    than the left (i.e.\ center) ones. The  margin on the bottom should be 
    larger than  on the top.   These rules  are automatically used  by the 
    KomaScript classes  (scrartcl, scrbook, scrrprt).  Do  not meddle with 
    them.                                                                  
    
    To   meddle  with   the   margins,   use  the   DIV   option  of   the 
    \verb+documentclass+  command.    Larger  values  result   in  smaller 
    margins.  If you  want to bind your thesis, make  sure to include some 
    space for  it using BCOR\@.   Ask the print shop  how much to  use for 
    your  specific binding  needs.  Never,  ever, use  the \verb+geometry+ 
    package.                                                               

    \section{Paragraphs are Empty Lines}
    \label{sec:paras}

    Paragraphs are separated  by empty lines in the  \LaTeX{} code.  Never 
    use  the double  backslash for  this purpose.   \LaTeX{} inserts  page 
    breaks  preferably between  paragraphs.  A  double backslash  lets you 
    stay in the  same paragraph, so \LaTeX{} does not  know where it would 
    be good  to insert page  breaks, figures, tables, algorithms,  and the 
    like.                                                                  

    Many  people find  the  indentation  of the  first  line in  paragraph 
    odd.  Look into  a favorite book or newspaper of  yours and you'll see 
    that this way of introducing  new paragraphs is used everywhere.  Note 
    that it  allows you---even on the  top of a new  page---to see whether 
    the line is the first of a new paragraph or the continuation of an old 
    paragraph.                                                             

    \section{Typesetting Math as Part of Your Text}

    Formulas are part of the text, try  to use them like nouns and be sure 
    to include  punctuation to  keep the text  flowing.  For  example, the 
    pythagorean theorem,                                                   
    \begin{align}
        a^2 + b^2 = c^2,
    \end{align}
    is typically taught at school.

    Other random notes:

    \begin{itemize}
        \item
    Note that  as in \cref{sec:paras},  before and after a  formula, empty 
    lines will  add space  and indentation, and  possibly page  breaks, or 
    figures.  So  if a formula  is part of a  paragraph, do not  add empty 
    lines around it in the code.                                           

        \item
    Number all equations.  It makes it  easier to refer to them for others,
    even if you do not refer to them in your own text.

        \item
    While we're at  it, do not use the  \verb+equation+ or \verb+eqnarray+ 
    environments, use \verb+align+ instead.

        \item
    Take care to always put all  math in the math-environment, for example 
    the  variable $x$.   Note  that  it is  typeset  differently from  the 
    letter~x.

        \item
    Use  single-letter  names  for  variables   in  math  (math  is  not  a
    programming language!).   Function names  may have  multiple characters,
    such as $\tanh(x)$.  Note that they are not typeset in italics.  
    They  usually have  their own  commands, and  you can  define your  own
    functions, such as  $\somefunc(x)$.  If you don't do  this, the kerning
    will  be broken,  since \LaTeX{}  assumes that  you're multiplying  the
    variables  $s$, $o$,  $m$, $e$,  $f$, $u$,  $n$, and  $c$ and  typesets
    accordingly,  compare  $somefunc$  (math  mode,  without  spaces!)  and
    \emph{somefunc} (italics).

    \end{itemize}

    \section{Referencing Sections, Figures, etc.}
    Use  the \verb+cleverref+  package to  refer  to other  places in  the 
    document,  e.g.  \cref{chap:Introduction}.    The  package  allows  to 
    specify the reference style globally.  Therefore, it helps to refer in 
    a consistent  way.  Otherwise you  may quickly have different  ways of 
    referencing,   e.g.  \emph{Figure},   \emph{figure},  \emph{Fig.}   and
    \emph{fig.}.   The package  also ensures  that there  is no  linebreak 
    before the number of the figure/chapter etc.                           

    \begin{figure}
        \centering
        \begin{subfigure}[b]{0.3\textwidth}
            \includegraphics[width=\textwidth]{images/cat.jpeg}
            \caption{Nyan cat}
            \label{fig:cat1}
        \end{subfigure}%
        ~ %add desired spacing between images, e. g. ~, \quad, \qquad, \hfill etc.
        %(or a blank line to force the subfigure onto a new line)
        \begin{subfigure}[b]{0.3\textwidth}
            \includegraphics[width=\textwidth]{images/cat.jpeg}
            \caption{Nyan cat}
            \label{fig:cat2} 
        \end{subfigure}
        ~ %add desired spacing between images, e. g. ~, \quad, \qquad, \hfill etc.
        %(or a blank line to force the subfigure onto a new line)
        \begin{subfigure}[b]{0.3\textwidth}
            \includegraphics[width=\textwidth]{images/cat.jpeg}
            \caption{Nyan cat}
            \label{fig:cat3}  
        \end{subfigure}
        \caption[Nyan cats]{Pictures of Nyan cats}\label{fig:cats}
    \end{figure}



    \section{Citing Other Works}
    There are two ways to cite, which depends on whether the authors are an
    active part of your sentence:
    \begin{itemize}[noitemsep]
        \item \citet{muster} showed that\ldots{}
        \item This is currently a hot research topic \citep{muster}
    \end{itemize}
    In  \LaTeX{}, these  different  citation styles  are  reflected in  the
    \texttt{\textbackslash citet}  and  \texttt{\textbackslash citep}  commands. \textbf{Never  use  the  plain
    \texttt{\textbackslash cite}!} -- it is outdated.
    This distinction  is even more important  when you use a  numeric style
    for referencing, where in the second version only the number of the 
    reference is printed.   In the first case, a  number-only citation does
    not make much sense.

    In  your  bibliography, you  should  define  common strings,  such  as 
    conference  names, using  the  \verb|@string|  syntax.  Otherwise  you 
    might  quickly end  up with  a  lot of  different names  for the  same 
    conference, which is confusing and  inconsistent.  (Please have a look 
    at the bibliography file.)                                             

    \section{Floats}
    \subsection{Figures}
    You should  allow \LaTeX{} to place  the figures where it  wants.  The 
    same  goes  for  tables.   This  is  called  a  \verb+float+,  as  the 
    figure/table floats around.   Look in a well-typeset  book, and you'll 
    notice that figures aren't placed randomly in the text, they're rather 
    always  on the  top  of  the page  if  at  all possible. \LaTeX{}  has 
    options to  let figures float to  \emph{top} (\verb+t+), \emph{bottom} 
    (\verb+b+),   \emph{page  of   floats}   (\verb+p+)  and   \emph{here} 
    (\verb+h+).  The  default, \verb+tbp+, is  ok for most users,  in any 
    case, never use \verb+h+.  You can use the FloatBarrier command if you 
    want to make sure that a figure is within a specific section.          

    All figures have  a caption from which you can  understand most of the 
    figure content  and its significance.   They are also  referenced from 
    the main text (e.g., see \cref{fig:cats}).  Do not force linebreaks in 
    captions.  Do not deviate from these rules, they are strict.           

    Figures  should  provide a  short  name  for  every figure  in  square 
    brackets: $[$~$]$. If  you do not define a short  name, the whole text 
    of the caption will be displayed in the table of figures, which is 
    usually too long.
    
    Note that \cref{fig:cats} also shows how to use subfigures.



    \subsection{Tables}
    Never  use vertical  lines in  tables.  See  the documentation  of the 
    booktabs  package  for  an  explanation,  or  look  in  your  favorite 
    scientific  book/magazine to  understand  that this  is a  commonplace 
    rule.                                                                  

    A basic  example, taken  from the booktabs  documentation, is  given in
    \cref{tab::ex}.

    \begin{table}
        \centering
        \begin{tabular}{@{}lrr@{}} 
            \toprule
            \multicolumn{2}{c}{Education}\\ \cmidrule{1-2}
            Major & Duration & Income (\euro)\\ 
            \midrule 
            CompSci & 2 & 12,75 \\ \addlinespace
            MST & 6 & 8,20 \\ \addlinespace
            VWL & 14 & 10,00\\ 
            \bottomrule
        \end{tabular}
        \caption[Table Example]{A basic example from the booktabs package.}
        \label{tab::ex}
    \end{table}

    \paragraph{The tabu package}
    The tabu package provides some useful tools for advanced tables.


    \subsection{Algorithms}
    There are several packages to typeset algorithms.
    We recommend the \verb+algorithm2e+ package.
    In \cref{alg:exp} a simple example from the algorithm2e documentation is given.

    \begin{algorithm}[t]
        \SetAlgoLined
        \KwData{this text}
        \KwResult{how to write algorithm with \LaTeX2e }
        initialization\;
        \While{not at end of this document}{
            read current\;
            \eIf{understand}{
                go to next section\;
                current section becomes this one\;
            }{
                go back to the beginning of current section\;
            }
        }
        \caption{How to write algorithms (Small example from the algorithm2e documentation)}
        \label{alg:exp}
    \end{algorithm}

    \section{Miscellaneous}

    \subsection{Enumerations}

    Short enumerations look ugly in  standard \LaTeX{}, since they include 
    too much  space between the  items.  A ``normal'' enumerate  list with 
    very short paragraphs:                                                 
    \begin{enumerate}
        \item foo bar baz
        \item foo bar baz
        \item foo bar baz
    \end{enumerate}

    A list using enumitem to compress the inter item spaces, looks much better for short items:
    \begin{enumerate}[noitemsep]
        \item foo bar baz
        \item foo bar baz
        \item foo bar baz
    \end{enumerate}

    \subsection{Front, Main and Back Matters}

    Titlepage, lists of figures, table of contents page etc.\ are numbered 
    in  roman letters.   The introduction  should start  with page  $1$ in 
    arabic  numerals.   The document  class  \verb+scrbook+  does this  by 
    default,  when \textbackslash  frontmatter, \textbackslash  mainmatter 
    and \textbackslash backmatter are placed properly.                     



    \section{Compiling The Document}
    \LaTeX{}  is  a bit  archaic  to  work  with,  since with  every  pass 
    through  your document,  some intermediate  information is  written to 
    files,  which are  then processed  by  other programs.   Keep in  mind 
    that  everytime a  new  pass through  your  document incorporates  new 
    information (e.g. from BibTeX), page breaks may change, which requires
    an  additional pass  through the  document. Thus, you  usually have  to
    compile the document in three steps:
    \begin{enumerate}
        \item Compile the document using xelatex: xelatex thesis
        \item Run biber  (or if  you do  not have  it, bibtex)  to set  up
              bibliography: biber thesis
        \item Run xelatex  twice more to incorporate  bibtex references and
              to get the typesetting right.
    \end{enumerate}

    If  you know  what  you're doing  you  can get  away  with less.   Some
    programs parse the  output of \LaTeX{} to  determine whether additional
    passes are necessary.  Most GUI programs fall in this category, as do
    \verb+rubber+ and \verb+latexmk+.

    \subsection{Issues With This Template}
    \begin{description}
        \item [My references can't be found] 
            The line including the \texttt{biblatex} package says it should
            use \texttt{biber} as a backend.  This is a reimplementation of
            BibTeX. If you do not have \texttt{biber}, or your IDE does not
            support it, consider switching the backend to \texttt{bibtex}.
    \end{description}

\end{document}
